%!TEX program = xelatex
% !TEX encoding = ISO-8859-1
\def\ano{2019}
\def\semestre{2}
\def\disciplina{\'Algebra 1}
\def\turma{C}
\def\numerolista{3}

\documentclass[12pt]{exam}

\usepackage{caption}
\usepackage{amssymb}
\usepackage{amsmath,amsfonts,amsthm,amstext}
\usepackage[brazil]{babel}
% \usepackage[latin1]{inputenc}
\usepackage{graphicx}
\graphicspath{{/home/jfreitas/GitHub_Repos/Algebra-1/Pictures/}{D:/OneDrive - unb.br/imagens-latex/}}
\usepackage{enumitem}
\usepackage{multicol}
\usepackage{answers}
\usepackage{tikz,ifthen}
\usetikzlibrary{lindenmayersystems}
\usetikzlibrary[shadings]
\Newassociation{solucao}{Solution}{ans}
\newtheorem{exercicio}{}

\setlength{\topmargin}{-1.0in}
\setlength{\oddsidemargin}{0in}
\setlength{\textheight}{10.1in}
\setlength{\textwidth}{6.5in}
\setlength{\baselineskip}{12mm}

\extraheadheight{0.7in}
\firstpageheadrule
\runningheadrule
\lhead{
        \begin{minipage}[c]{1.7cm}
        \includegraphics[width=1.7cm]{unb.pdf}
        \end{minipage}%
        \hspace{0pt}
        \begin{minipage}[c]{4in}
          {Universidade de Bras{\'\i}lia} --
          {Departamento de Matem{\'a}tica}
\end{minipage}
\vspace*{-0.8cm}
}
% \chead{Universidade de Bras{\'\i}lia - Departamento de Matem{\'a}tica}
% \rhead{}
% \vspace*{-2cm}

\extrafootheight{.5in}
\footrule
\lfoot{\disciplina\ - \semestre$^o$/\ano}
\cfoot{}
\rfoot{P\'agina \thepage\ de \numpages}

\newcounter{exercicios}
\renewcommand{\theexercicios}{\arabic{exercicios}}

\newenvironment{questao}[1]{
\refstepcounter{exercicios}
\ifx&#1&
\else
   \label{#1}
\fi
\noindent\textbf{Exerc{\'\i}cio {\theexercicios}:}
}

\newcommand{\resp}[1]{
\noindent{\bf Exerc{\'\i}cio #1: }}
\newcommand{\nsub}{\varsubsetneq}
\newcommand{\vaz}{\emptyset}
\newcommand{\im}{{\rm Im\,}}
\newcommand{\sub}{\subseteq}
\newcommand{\n}{\mathbb{N}}
\newcommand{\z}{\mathbb{Z}}
\newcommand{\rac}{\mathbb{Q}}
\newcommand{\real}{\mathbb{R}}
\newcommand{\complex}{\mathbb{C}}
\newcommand{\cp}[1]{\mathbb{#1}}
\newcommand{\ch}{\mbox{\textrm{car\,}}\nobreak}
\newcommand{\vesp}[1]{\vspace{ #1  cm}}
\newcommand{\compcent}[1]{\vcenter{\hbox{$#1\circ$}}}
\newcommand{\comp}{\mathbin{\mathchoice
{\compcent\scriptstyle}{\compcent\scriptstyle}
{\compcent\scriptscriptstyle}{\compcent\scriptscriptstyle}}}

\begin{document}

\begin{center}
{\Large\bf \disciplina\ - Turma \turma\ -- \semestre$^{o}$/\ano} \\ \vspace{9pt} {\large\bf
  $\numerolista^{\underline{a}}$ Lista de Exerc{\'\i}cios -- Operações em $\z_m$}\\ \vspace{9pt} Prof. Jos{\'e} Ant{\^o}nio O. Freitas
\end{center}
\hrule

\vspace{.6cm}

\questao{} Construa a tabela de soma, $\oplus$, e multiplicação, $\otimes$, nos seguintes conjuntos:
\begin{multicols}{3}
	\begin{enumerate}[label=({\alph*})]
		\item $\z_5$
		\item $\z_6$
		\item $\z_8$
		\item $\z_9$
		\item $\z_{11}$
		\item $\z_{12}$
	\end{enumerate}
\end{multicols}

\vspace{.3cm}

\questao{} Considere os seguintes subconjuntos $G$ de $\z_{12}$:
\begin{multicols}{2}
	\begin{enumerate}[label=({\alph*})]
		\item $G=\{\overline{1},\overline{11}\}$;

		\item $G=\{\overline{0},\overline{4},\overline{8}\}$;

		\item $G=\{\overline{0},\overline{2},\overline{4},\overline{6},\overline{8},\overline{10}\}$
		\item $G=\{\overline{1}, \overline{3},\overline{5},\overline{7},\overline{9},\overline{11}\}$.
	\end{enumerate}
\end{multicols}

Para quais desses conjuntos vale as duas propriedades seguintes:
\begin{enumerate}[label=({\roman*})]
	\item $x \oplus y \in G$, para todos $x$, $y \in G$.
	\item Para todo $x \in G$, existe $y \in G$ tal que $x \oplus y = \overline{0}$.
\end{enumerate}

\vspace{.3cm}

\questao{} Considere os seguintes subconjuntos $G$ de $\z_{13}$:
\begin{multicols}{2}
	\begin{enumerate}[label=({\alph*})]
		\item $G=\{\overline{1},\overline{12}\}$;

		\item $G=\{\overline{1},\overline{5},\overline{8},\overline{12}\}$;

		\item $G=\{\overline{1},\overline{2},\overline{3},\overline{4}, \overline{5},\overline{6},\overline{7},
		 \overline{8},\overline{9},\overline{10},\overline{11},\overline{12}\}$
		\item $G=\{\overline{1}, \overline{3},\overline{5},\overline{7},\overline{9},\overline{11}\}$.
	\end{enumerate}
\end{multicols}

Para quais desses conjuntos vale as duas propriedades seguintes:
\begin{enumerate}[label=({\roman*})]
	\item $x \otimes y \in G$, para todos $x$, $y \in G$.
	\item Para todo $x \in G$, existe $y \in G$ tal que $x \otimes y = \overline{1}$.
\end{enumerate}

\vspace{.3cm}

\questao{} Resolva as seguintes equa\c{c}\~oes:
\begin{enumerate}[label={\alph*})]
	\item $x \oplus \overline{2} = \overline{0}$ no conjunto $\z_4$;
	\item $x^2 \oplus \overline{2}x = \overline{2}$ no conjunto $\z_5$;
	\item $x^3 \oplus \overline{2}x = \overline{3}$ no conjunto $\z_7$.
\end{enumerate}

\newpage

\questao{} Resolva o sistema de equa\c{c}\~oes
  \[
    \begin{cases}
      x \oplus y = \overline{0}\\
      \overline{2}x \oplus y = \overline{2}
    \end{cases}
  \]
nos conjuntos:
\begin{multicols}{2}
	\begin{enumerate}[label=({\alph*})]
		\item $\z_4$
		\item $\z_6$
		\item $\z_3$
		\item $\z_5$
	\end{enumerate}
\end{multicols}

\vspace{.3cm}

\questao{} Resolva a equa\c{c}\~ao $x^2 \oplus \overline{2}x \oplus \overline{1} = \overline{0}$ nos conjuntos $\z_7$ e $\z_{11}$, caso ela tenha ra{\'\i}zes.

\vspace{.3cm}

\questao{} Para quais valores de $\overline{c}$ a equa\c{c}\~ao $\overline{2}\otimes x = \overline{c}$ tem solu\c{c}\~ao no conjunto $\z_5$.

\vspace{.3cm}

\questao{} Mostre que a rela{\c c}{\~a}o de congru{\^e}ncia m{\'o}dulo $m$ satisfaz as seguintes propriedades:
\begin{enumerate}[label={\roman*})]
	\item $a_{1}\equiv b_{1}\pmod{m}$ se, e somente se, $a_{1}-b_{1}\equiv 0\pmod{m}$.
	\item Se $a_{1}\equiv b_{1}\pmod{m}$ e $a_{2}\equiv b_{2}\pmod{m}$, ent{\~a}o $a_{1}+a_{2}\equiv b_{1}+b_{2}\pmod{m}$.
	\item Se $a_{1}\equiv b_{2}\pmod{m}$ e $a_{2}\equiv b_{2}\pmod{m}$, ent{\~a}o $a_{1}a_{2}\equiv b_{1}b_{2}\pmod{m}$.\label{item_provado}
	\item Se $a\equiv b\pmod{m}$, ent{\~a}o $ax\equiv bx\pmod{m}$, para todo $x \in \z$.
	\item Vale a lei do cancelamento: se $d \in \z$ e $mdc(d,m) = 1$ ent{\~a}o $ad \equiv bd \pmod{m}$ implica $a\equiv b \pmod{m}$.
\end{enumerate}
\end{document}